\documentclass[11pt]{article}
\usepackage[letterpaper,margin=1in]{geometry}
\usepackage[T1]{fontenc}
\usepackage[utf8]{inputenc}
\usepackage{lmodern,microtype}
\usepackage{amsmath,amssymb,mathtools,bm,slashed}
\usepackage{xcolor,enumitem,booktabs,array,needspace}
\usepackage[most]{tcolorbox}
\usepackage{fancyhdr,titlesec,parskip,hyperref}
\definecolor{qftblue}{HTML}{245A8D}
\definecolor{qftgreen}{HTML}{2F6B4F}
\definecolor{qftgray}{HTML}{F4F6F8}
\hypersetup{colorlinks=true,linkcolor=qftblue,urlcolor=qftblue,pdfauthor={Nima Lashkari}}
\pagestyle{fancy}\fancyhf{}
\lhead{PHYS 662: Quantum Field Theory I}\cfoot{\thepage}
\setlength{\headheight}{14pt}
\titleformat{\section}{\Large\bfseries\color{qftblue}}{\thesection.}{0.6em}{}
\titleformat{\subsection}{\large\bfseries\color{qftgreen}}{\thesubsection}{0.6em}{}
\setlist[enumerate]{leftmargin=2em,itemsep=0.25em,topsep=0.35em}
\setlist[itemize]{leftmargin=2em,itemsep=0.2em,topsep=0.35em}
\newcommand{\dd}{\mathrm d}
\newcommand{\ii}{\mathrm i}
\newcommand{\eps}{\varepsilon}
\newcommand{\tr}{\operatorname{tr}}
\newcommand{\Tr}{\operatorname{Tr}}
\newcommand{\diag}{\operatorname{diag}}
\newcommand{\Lag}{\mathcal L}
\newcommand{\D}{\mathcal D}
\newcommand{\vac}{\lvert0\rangle}
\newcommand{\vbra}{\langle0\rvert}
\newcommand{\Tord}{\mathrm T}
\newcommand{\ket}[1]{\lvert#1\rangle}
\newcommand{\bra}[1]{\langle#1\rvert}
\newtcolorbox{keyidea}[1][]{enhanced,breakable,colback=qftgray,colframe=qftblue,
 boxrule=0.7pt,arc=1.5mm,title={#1},fonttitle=\bfseries,coltitle=white}
\newtcolorbox{remarkbox}[1][]{enhanced,breakable,colback=white,colframe=qftgreen,
 boxrule=0.6pt,arc=1.5mm,title={#1},fonttitle=\bfseries,coltitle=white}
\newcommand{\intp}{\int\!\frac{\dd^4p}{(2\pi)^4}}
\newcommand{\MSbar}{\overline{\mathrm{MS}}}
\newcommand{\module}[1]{\clearpage\section{#1}}
\newcommand{\braket}[2]{\langle#1\mid#2\rangle}
\newcommand{\dPi}{\dd\Pi}
\newcommand{\cM}{\mathcal M}
\newcommand{\cH}{\mathcal H}
\newcommand{\cS}{\mathcal S}
\newcommand{\Zp}{Z_{\rm pole}}
\newcommand{\Kallen}{\mathcal K}
\newcommand{\lecturefront}[4]{%
\rhead{Lecture #1}\hypersetup{pdftitle={PHYS 662 Lecture #1: #2}}
\begin{document}\pagenumbering{gobble}
\begin{titlepage}\centering\vspace*{1.05in}
{\scshape\Large PHYS 662 --- Quantum Field Theory I\par}
\vspace{0.35in}{\Huge\bfseries\color{qftblue}Lecture #1\par}
\vspace{0.2in}{\LARGE\bfseries #2\par}
\vspace{0.2in}{\large\color{qftgreen}#3\par}
\vspace{0.6in}{\large Nima Lashkari\par}{\large Fall 2026\par}
\vfill\begin{minipage}{0.86\textwidth}\small #4\par\medskip
Course page: \url{https://nimalashkari.com/quantum-field-theory-2026-part-i/}
\end{minipage}\vspace{0.5in}\end{titlepage}
\pagenumbering{roman}\setcounter{tocdepth}{1}\tableofcontents
\vfill\begin{remarkbox}[Place in the course]
Lecture 26 constructs the physical photon Hilbert space and its propagators.
Lecture 27 develops QED amplitudes. Lecture 28 computes the photon and
fermion self-energies. Lecture 29 relates the vertex to charge
renormalization and the magnetic moment. Homework 9 closes this sequence.
\end{remarkbox}
\small\textbf{Reading and conventions.} The uploaded course slides are the
primary starting point. D.~Tong,
\href{https://www.damtp.cam.ac.uk/user/tong/qft/qfthtml/S6.html}{\textit{Quantum Field Theory}, Chapter 6},
is a companion for photon quantization and QED. For the vertex calculation,
see V.~Kaplunovsky,
\href{https://web2.ph.utexas.edu/~vadim/Classes/2020f/vertex.pdf}{\textit{QED Vertex Correction}},
especially the parameter integral for the Pauli form factor.
For further textbook study, see Peskin--Schroeder, \textit{An Introduction
to QFT}, Chapters 5--7 and 10.
All formulas here use the mostly-minus metric and propagators that include
\(\ii\). Dimensional regularization uses \(d=4-2\eps\), which differs from
some of the slides. The accompanying editorial memo records corrections
and the scope of the supplementary reading.
\normalsize\clearpage\pagenumbering{arabic}}
\lecturefront{26}{Quantization of the Electromagnetic Field}
{Constraints, physical photons, and covariant propagators}
{The electromagnetic potential has four components, but a photon has only
two physical polarizations. We follow this mismatch from the classical
constraint equations through canonical quantization and the path integral,
and explain why gauge-dependent intermediate expressions give consistent
physical amplitudes.}

\section{Maxwell theory is a constrained system}
We use \(\eta_{\mu\nu}=\diag(1,-1,-1,-1)\), \(\hbar=c=1\), and
\(A^\mu=(A^0,\bm A)\). With a prescribed conserved current
\(j^\mu=(\rho,\bm J)\),
\begin{equation}
 \Lag=-\frac14 F_{\mu\nu}F^{\mu\nu}-j^\mu A_\mu,
 \qquad F_{\mu\nu}=\partial_\mu A_\nu-\partial_\nu A_\mu.
\end{equation}
Thus \(\bm E=-\dot{\bm A}-\bm\nabla A^0\),
\(\bm B=\bm\nabla\times\bm A\), and the kinetic term is
\((\bm E^2-\bm B^2)/2\). Varying \(A_\nu\) gives
\begin{equation}
 \partial_\mu F^{\mu\nu}=j^\nu.
 \label{eq:maxwell}
\end{equation}
Its divergence requires \(\partial_\nu j^\nu=0\). The same condition makes
the source action invariant under \(A_\mu\mapsto A_\mu+\partial_\mu\chi\),
up to a boundary term. Gauge-related potentials describe the same field
strength; the redundancy is local in spacetime.

The Lagrangian contains no \(\dot A_0\). Using the covariant components
\(A_\mu\) as coordinates, the momenta are
\begin{equation}
 \pi^\mu=F^{\mu0},\qquad \pi^0=0,\qquad \pi^i=E^i.
\end{equation}
The relation \(\pi^0=0\) is a primary constraint. The \(\nu=0\) equation
in \eqref{eq:maxwell} is Gauss's law,
\begin{equation}
 \bm\nabla\cdot\bm E=\rho.
\end{equation}
It constrains the initial data rather than furnishing a second-order
evolution equation for \(A^0\). In a Hamiltonian description, preserving
the primary constraint in time produces this secondary constraint.

\begin{keyidea}[Why four oscillators would be too many]
The two first-class constraints and their gauge freedom remove four
phase-space coordinates from the eight associated with \(A_\mu,\pi^\mu\).
The remaining four phase-space coordinates describe two propagating
configuration degrees of freedom. This count assumes the usual boundary
conditions and omits possible global zero modes.
\end{keyidea}
Setting \(A^0=0\) before varying the action would hide Gauss's law.
Temporal gauge can be used, but Gauss's law must then be imposed separately.
We instead begin with Coulomb gauge, where its solution is explicit.

\module{Coulomb gauge and the instantaneous interaction}
Choose \(\bm\nabla\cdot\bm A=0\), with fields and gauge transformations
vanishing suitably at spatial infinity. Write \(\bm A=\bm A_T\).
Gauss's law becomes a Poisson equation,
\begin{equation}
 -\nabla^2 A^0(t,\bm x)=\rho(t,\bm x),\qquad
 A^0(t,\bm x)=\int\dd^3y\,
 \frac{\rho(t,\bm y)}{4\pi|\bm x-\bm y|}.
\end{equation}
The cross term \(\int\dot{\bm A}_T\cdot\bm\nabla A^0\) vanishes by
integration by parts. Substituting the solution into the action gives
\begin{align}
 L_{\rm red}={}&\int\dd^3x\left[
 \frac12\dot{\bm A}_T^2-\frac12(\bm\nabla\times\bm A_T)^2
 +\bm J_T\cdot\bm A_T\right]-E_C,\\
 E_C={}&\frac12\int\dd^3x\,\dd^3y\,
 \frac{\rho(t,\bm x)\rho(t,\bm y)}{4\pi|\bm x-\bm y|}.
 \label{eq:coulombenergy}
\end{align}
The factor \(1/2\) avoids counting each pair twice. For two point sources,
the cross term is \(q_1q_2/(4\pi r)\); individual point-charge self-energies
require regulation. The sign in the reduced Lagrangian is \(-E_C\), while
the Hamiltonian contains \(+E_C\).

The transverse projection is
\begin{equation}
 P_T^{ij}(\bm k)=\delta^{ij}-\frac{k^ik^j}{|\bm k|^2},\qquad
 J_T^i=P_T^{ij}J^j.
\end{equation}
For the physical spatial coordinates \(A_T^i\), the conjugate momentum is
\(\Pi_T^i=\dot A_T^i=-E_T^i\). This sign differs from \(\pi^i\) above
because \(A_i=-A^i\). The canonical commutator on the reduced space is
\begin{equation}
 [A_T^i(t,\bm x),\Pi_T^j(t,\bm y)]
 =\ii\delta_T^{ij}(\bm x-\bm y),\quad
 \delta_T^{ij}(\bm r)=\int\frac{\dd^3k}{(2\pi)^3}
 P_T^{ij}(\bm k)e^{\ii\bm k\cdot\bm r}.
\end{equation}
An unrestricted \(\delta^{ij}\delta^{(3)}\) would contradict
\(\partial_iA_T^i=0\). Projection therefore changes the commutator as well
as the equations of motion.

\begin{remarkbox}[Instantaneous does not mean acausal]
The Coulomb term is an interaction of constrained variables, not a new
propagating mode. The transverse field supplies the remaining pieces of
electromagnetic response. Their combination gives causal gauge-invariant
field strengths; an isolated gauge-dependent potential need not have a
manifestly causal representation.
\end{remarkbox}

\module{Two transverse photon oscillators}
For the free field choose spatial polarization vectors obeying
\begin{equation}
 \bm k\cdot\bm\epsilon_\lambda=0,\quad
 \bm\epsilon_\lambda^*\cdot\bm\epsilon_{\lambda'}=\delta_{\lambda\lambda'},
 \quad \sum_{\lambda=1}^2\epsilon_\lambda^i\epsilon_\lambda^{j*}
 =P_T^{ij}(\bm k).
\end{equation}
Let \(\omega_k=|\bm k|\) and
\(\dd\Pi_k=\dd^3k/[(2\pi)^3 2\omega_k]\). A normalization consistent
with our earlier particle states is
\begin{equation}
 A_T^i(x)=\sum_{\lambda=1}^2\int\dd\Pi_k
 \left[\epsilon_\lambda^i(\bm k)a_\lambda(\bm k)e^{-\ii kx}
 +\epsilon_\lambda^{i*}(\bm k)a_\lambda^\dagger(\bm k)e^{\ii kx}\right],
 \quad k^0=\omega_k,
\end{equation}
with
\begin{equation}
 [a_\lambda(\bm k),a_{\lambda'}^\dagger(\bm q)]
 =2\omega_k(2\pi)^3\delta_{\lambda\lambda'}\delta^{(3)}(\bm k-\bm q).
\end{equation}
Differentiating the mode expansion and using completeness reproduces the
transverse canonical commutator. This is a useful normalization check:
one power of \(2\omega_k\) comes from the commutator and the other from
the time derivatives and the two frequency contributions.

Substitution into the free energy, followed by normal ordering, yields
\begin{equation}
 :H_\gamma:=\sum_\lambda\int\dd\Pi_k\,
 \omega_k a_\lambda^\dagger(\bm k)a_\lambda(\bm k),\qquad
 :\bm P_\gamma:=\sum_\lambda\int\dd\Pi_k\,
 \bm k a_\lambda^\dagger(\bm k)a_\lambda(\bm k).
\end{equation}
The \(aa\) and \(a^\dagger a^\dagger\) terms cancel between electric and
magnetic energy. A state \(\ket{k,\lambda}=a_\lambda^\dagger(\bm k)\vac\)
has positive norm, energy \(\omega_k\), and the desired covariant
normalization. There is no longitudinal photon oscillator in this construction.

For momentum in the \(z\) direction, the two linear polarizations can be
combined into circular polarizations
\(\bm\epsilon_\pm=(\bm e_x\pm\ii\bm e_y)/\sqrt2\).
They transform by opposite phases under rotations about \(\bm k\), and
give the two helicities \(\pm1\). The association of the labels \(+,-\)
with a chosen phase convention is secondary; the two rotation eigenvalues
are physical. A massless photon has no rest frame and no third spin
projection obtained by boosting a rest-frame vector.

\begin{keyidea}[What has been quantized]
The free photon Fock space is the symmetric Fock space over the two
helicity states. The Coulomb interaction will reappear when charged
matter is present, but it does not add an independent polarization.
\end{keyidea}

\module{Polarization vectors and gauge equivalence}
A covariant polarization representative satisfies
\begin{equation}
 k^2=0,\qquad k\cdot\epsilon=0,\qquad
 \epsilon^\mu\sim\epsilon^\mu+c\,k^\mu.
 \label{eq:polquotient}
\end{equation}
The last equivalence follows because adding a plane-wave gradient to
\(A_\mu\) leaves \(F_{\mu\nu}\) unchanged. The space orthogonal to a null
vector is three-dimensional, but it contains that null vector itself.
The quotient in \eqref{eq:polquotient} is therefore two-dimensional.
This is the covariant version of our transverse oscillator count.

To choose representatives, introduce a vector \(n\) with \(k\cdot n\ne0\)
and impose \(n\cdot\epsilon=0\). With
\(\epsilon_\lambda\cdot\epsilon_{\lambda'}^*=-\delta_{\lambda\lambda'}\),
\begin{equation}
 \sum_{\lambda=1}^2\epsilon_\lambda^\mu\epsilon_\lambda^{\nu*}
 =-\eta^{\mu\nu}
 +\frac{k^\mu n^\nu+n^\mu k^\nu}{k\cdot n}
 -\frac{n^2 k^\mu k^\nu}{(k\cdot n)^2}.
 \label{eq:polsum}
\end{equation}
Contracting either index with \(k\) or \(n\) gives zero. Setting
\(n=(1,\bm0)\) makes the time components vanish and reduces the spatial
components to \(P_T^{ij}\).

Suppose an amplitude with one external photon is
\(\cM=\epsilon_\mu\mathcal A^\mu\). Independence of the representative
requires the on-shell Ward identity
\begin{equation}
 k_\mu\mathcal A^\mu=0.
\end{equation}
Only after this identity is established may we replace the physical
polarization sum by \(-\eta^{\mu\nu}\) in the squared amplitude:
all other terms in \eqref{eq:polsum} then vanish. This shortcut sums two
physical states by an algebraic identity; it does not assert the existence
of four physical photon polarizations.

\begin{remarkbox}[An off-shell projector is a different object]
For \(k^2\ne0\), \(\eta^{\mu\nu}-k^\mu k^\nu/k^2\) has rank three.
It is useful for inverting quadratic operators and decomposing loop
corrections. It is singular at a null momentum and is not the sum over
on-shell photon polarizations.
\end{remarkbox}
A Lorentz transformation of a Coulomb-gauge polarization need not stay in
Coulomb gauge. A compensating gauge transformation restores the chosen
representative, while its physical helicity state transforms consistently.

\module{Gauge fixing in the path integral}
The Maxwell quadratic action is
\begin{equation}
 S_0=\frac12\int\dd^4x\, A_\mu
 (\eta^{\mu\nu}\Box-\partial^\mu\partial^\nu)A_\nu.
\end{equation}
The kinetic operator annihilates every gradient \(\partial_\nu\chi\).
Consequently its determinant vanishes and it has no inverse on the full
space of potentials. The naive Gaussian functional integral counts each
gauge orbit infinitely many times.

Choose the gauge condition \(G[A]=\partial\cdot A=f\). Along an orbit,
\(G[A+\partial\chi]=G[A]+\Box\chi\). Formally inserting
\begin{equation}
 1=\det(\Box)\int\D\chi\,
 \delta\bigl(\partial\cdot A+\Box\chi-f\bigr)
\end{equation}
selects one representative, with the gauge-group volume subsequently
removed. Boundary conditions must eliminate residual transformations
with \(\Box\chi=0\), or their zero modes must be treated separately.
This qualification is part of the definition of the determinant.

Average over \(f\) with weight
\(\exp[-\ii\int f^2/(2\xi)]\). The result is the gauge-fixed action
\begin{equation}
 \Lag_\xi=-\frac14F^2-\frac1{2\xi}(\partial\cdot A)^2.
 \label{eq:gaugefix}
\end{equation}
The determinant \(\det\Box\) is independent of \(A\). In normalized Abelian
correlators it cancels between numerator and denominator. If represented
by Grassmann ghost fields, those fields are free and decouple. There are
therefore no interacting Faddeev--Popov ghost loops in ordinary QED.
This statement will change for non-Abelian gauge fields.

\begin{keyidea}[A gauge choice defines intermediate Green functions]
The parameter \(\xi\) changes \(\langle\Tord A_\mu A_\nu\rangle\).
Gauge-invariant quantities are independent of it when the calculation
and regulator respect the Ward identities. Gauge fixing makes a redundant
Gaussian invertible; it does not give the photon a mass.
\end{keyidea}
The Landau limit is \(\xi\to0\), which enforces the Lorenz condition
\(\partial\cdot A=0\). Feynman gauge is \(\xi=1\). The phrase
``Lorenz gauge'' alone should not be used to identify a unique member of
this covariant gauge family.

\module{Inverting the covariant quadratic operator}
For \(k^2\ne0\), define
\(P_L^{\mu\nu}=k^\mu k^\nu/k^2\) and
\(P_T^{\mu\nu}=\eta^{\mu\nu}-P_L^{\mu\nu}\).
They are orthogonal projectors as mixed-index matrices. Fourier
transformation of \eqref{eq:gaugefix} gives
\begin{equation}
 K^{\mu\nu}(k)=-k^2P_T^{\mu\nu}-\frac{k^2}{\xi}P_L^{\mu\nu}.
\end{equation}
The Gaussian contraction is \(\ii K^{-1}\), with the Feynman boundary
prescription. Thus, with that prescription understood at each pole,
\begin{equation}
 D_{\mu\nu}^{(\xi)}(k)
 =\frac{-\ii}{k^2+\ii0}
 \left[\eta_{\mu\nu}-(1-\xi)\frac{k_\mu k_\nu}{k^2}\right].
 \label{eq:photonprop}
\end{equation}
The longitudinal double pole is itself a distribution defined by the
same prescription, not a reason to discard the term. Away from the pole,
\(k^\mu D_{\mu\nu}^{(\xi)}=-\ii\xi k_\nu/k^2\).
Only Landau gauge is transverse as a propagator.

In particular,
\begin{equation}
 D_{\mu\nu}^{(1)}(k)=\frac{-\ii\eta_{\mu\nu}}{k^2+\ii0}.
\end{equation}
The minus sign relative to the scalar propagator follows from the
Maxwell kinetic operator with our metric. Its spatial transverse
components have the same positive oscillator contraction as two free
massless scalar degrees of freedom.

For comparison, the Coulomb-gauge contractions are
\begin{equation}
 D_C^{00}(k)=\frac{\ii}{|\bm k|^2},\quad D_C^{0i}=0,\quad
 D_C^{ij}(k)=\frac{\ii P_T^{ij}(\bm k)}{(k^0)^2-|\bm k|^2+\ii0}.
 \label{eq:coulombprop}
\end{equation}
The \(00\) component has no energy pole: Fourier transformation in
\(k^0\) produces \(\delta(t-t')\), the instantaneous Coulomb kernel.
The spatial propagator contains the two radiation modes.

\begin{remarkbox}[Two different uses of transversality]
Physical external photons have transverse polarization representatives.
The photon self-energy will be transverse in every covariant gauge.
Neither statement implies that the full gauge-fixed photon propagator is
transverse for arbitrary \(\xi\). Keeping these objects distinct prevents
sign and gauge errors in the Dyson series of Lecture 28.
\end{remarkbox}

\module{Covariant quantization and the physical space}
Covariant quantization retains four components of \(A^\mu\) temporarily.
In Feynman gauge the free mode algebra can be written
\begin{equation}
 [a^\mu(\bm k),a^{\nu\dagger}(\bm q)]
 =-\eta^{\mu\nu}2\omega_k(2\pi)^3\delta^{(3)}(\bm k-\bm q).
\end{equation}
The timelike oscillator has a negative one-particle norm. This enlarged
space is an auxiliary indefinite-metric space, not the physical photon
Hilbert space. The sign of one propagator component alone would not be
an adequate definition of the physical space.

Imposing \(\partial\cdot A=0\) as a strong operator equation conflicts
with the covariant canonical algebra. Even the condition
\((\partial\cdot A)\ket{\Psi}=0\) is too strong: its creation part does
not annihilate the vacuum. Instead impose the Gupta--Bleuler condition
\begin{equation}
 (\partial\cdot A)^{(+)}\ket{\Psi}=0,
 \quad\hbox{equivalently}\quad k_\mu a^\mu(\bm k)\ket{\Psi}=0.
\end{equation}
Here \((+)\) denotes the positive-frequency, annihilation part.
For momentum along \(z\) this is proportional to
\((a^0-a^3)\ket{\Psi}=0\).

For example, \((a^{0\dagger}-a^{3\dagger})\vac\) obeys the condition
and has zero norm: the negative timelike and positive longitudinal
contributions cancel. The constrained space is positive semidefinite,
with such null states still present. The actual Hilbert space is
\begin{equation}
 \cH_{\rm phys}=
 \frac{\ker(\partial\cdot A)^{(+)}}{\text{null states}}.
\end{equation}
Its positive-norm excitations are precisely the two transverse photons.
Gauge-invariant observables act consistently on this quotient, and matrix
elements of \(\partial\cdot A\) between physical states vanish.

\begin{keyidea}[The two quantizations agree]
Coulomb gauge solves the constraints before quantizing and displays
positivity directly. Covariant quantization makes Lorentz covariance
simple, then removes the auxiliary states by a subsidiary condition and
a quotient. Both describe the same free physical photon states.
\end{keyidea}
With interactions, current conservation and the Ward identities ensure
that unphysical polarizations decouple from physical amplitudes.
For Abelian QED this provides the bridge from the free construction to
the scattering rules; it does not license arbitrary off-shell
polarization replacements.

\module{Conserved sources and the bridge to QED}
A simple check of gauge independence is exchange between conserved
currents. If \(k_\mu J^\mu=k_\mu J'^\mu=0\), then
\begin{equation}
 J_\mu D_{(\xi)}^{\mu\nu}J'_\nu
 =\frac{-\ii J\cdot J'}{k^2+\ii0},
\end{equation}
independent of \(\xi\). This cancellation uses current conservation,
not regularity of longitudinal terms at the photon pole.

The same identity follows less manifestly in Coulomb gauge. Let
\(K=|\bm k|\), and decompose the spatial current into its transverse
part and its component parallel to \(\bm k\). Conservation gives
\(\bm k\cdot\bm J=k^0J^0\). Away from poles,
\begin{align}
 J_\mu D_C^{\mu\nu}J'_\nu
 &=\frac{\ii J^0J'^0}{K^2}
   +\frac{\ii\bm J_T\cdot\bm J'_T}{k^2},\\
 J\cdot J'&=-\frac{k^2}{K^2}J^0J'^0
             -\bm J_T\cdot\bm J'_T.
\end{align}
The two expressions reproduce the covariant result, with the appropriate
boundary prescription. The instantaneous and radiation contributions
have combined into a Lorentz-covariant answer.

Now make the current dynamical by introducing a Dirac field. Denote its
\emph{signed} charge by \(e\), negative for the electron, and set
\begin{equation}
 D_\mu\psi=(\partial_\mu+\ii e A_\mu)\psi,\qquad
 \Lag_{\rm QED}=-\frac14F^2+\bar\psi(\ii\gamma^\mu D_\mu-m)\psi.
\end{equation}
The local transformations
\begin{equation}
 A_\mu\mapsto A_\mu+\partial_\mu\chi,\qquad
 \psi\mapsto e^{-\ii e\chi}\psi
\end{equation}
make \(D_\mu\psi\) transform like \(\psi\). The interaction is
\(-e\bar\psi\gamma^\mu\psi A_\mu\), and
\(j^\mu=e\bar\psi\gamma^\mu\psi\).

\begin{keyidea}[Ready for amplitudes]
An internal photon uses the gauge-fixed propagator. An incoming physical
photon contributes \(\epsilon_\mu(k,\lambda)\); an outgoing photon
contributes \(\epsilon_\mu^*(k,\lambda)\). Matter couples through
\(-\ii e\gamma^\mu\). Physical polarization sums may be simplified only
after checking the Ward identity for the complete amplitude.
\end{keyidea}
Lecture 27 applies these rules to exchange, annihilation, and Compton
scattering. Homework 9, Problem 1, tests their common origin in the
constraint structure of Maxwell theory.
\end{document}
